\documentclass[11pt,a4paper]{article}
\usepackage[margin=16mm,top=49mm,bottom=19mm,headheight=34mm,headsep=5mm]{geometry}
\usepackage{fontspec}
\setmainfont{Times New Roman}
\usepackage{amsmath,amssymb}
\usepackage{graphicx,array,longtable}
\usepackage[table]{xcolor}
\usepackage{fancyhdr}
\usepackage[hidelinks]{hyperref}
\hypersetup{pdftitle={Program Remedial dan Pengayaan Matematika Fase F XI - Barisan dan Deret},pdfauthor={Imran, S.Si.}}
\setlength{\parindent}{0pt}
\setlength{\LTpre}{3pt}\setlength{\LTpost}{4pt}
\setlength{\emergencystretch}{2em}
\linespread{1.08}
\pagestyle{fancy}\fancyhf{}
\renewcommand{\headrulewidth}{0pt}
\fancyhead[C]{\begin{minipage}{\textwidth}
\noindent\begin{minipage}[c]{23mm}
\centering\includegraphics[width=20mm,height=23mm,keepaspectratio]{logo.png}
\end{minipage}%
\begin{minipage}[c]{132mm}\centering
{\fontsize{10.2}{12}\selectfont\bfseries KEMENTERIAN AGAMA REPUBLIK INDONESIA\\
KANTOR KEMENTERIAN AGAMA KABUPATEN GOWA}\\[2pt]
{\fontsize{13}{15}\selectfont\bfseries MAN INSAN CENDEKIA GOWA}\\[2pt]
{\fontsize{8.3}{10.5}\selectfont
\mbox{Jl. Pendidikan Desa Belapunrangan, Kec. Parangloe, Gowa, Sulawesi Selatan 92173}\\
\mbox{Telp. 085173279846 $\cdot$ Pos-el: insancendekiagowa@gmail.com $\cdot$ Laman: icgowa.sch.id}}
\end{minipage}%
\begin{minipage}[c]{23mm}\strut\end{minipage}
\par\vspace{2mm}\hrule height1.2pt\vspace{1.2pt}\hrule height.4pt
\end{minipage}}
\fancyfoot[L]{\fontsize{8}{10}\selectfont Program Remedial dan Pengayaan Matematika $|$ Imran, S.Si.}
\fancyfoot[R]{\fontsize{8}{10}\selectfont\thepage}
\begin{document}
\begin{center}\bfseries\large PROGRAM REMEDIAL DAN PENGAYAAN\\[2pt] MATEMATIKA FASE F KELAS XI\\[2pt] MATERI BARISAN DAN DERET\end{center}\medskip
\begin{center}TAHUN PELAJARAN 2026/2027\end{center}
\vspace{19.84251969pt}
{\small\setlength{\tabcolsep}{2mm}\renewcommand{\arraystretch}{1.35}\begin{longtable}{|>{\raggedright\arraybackslash}p{49.324mm}|>{\raggedright\arraybackslash}p{120.226mm}|}\hline
Satuan pendidikan & MAN Insan Cendekia Gowa\\ \hline
Mata pelajaran & Matematika\\ \hline
Fase / kelas & F / XI\\ \hline
Elemen & Bilangan\\ \hline
Materi & Barisan dan Deret\\ \hline
Guru penyusun & Imran, S.Si.\\ \hline
NIP & 199309082019031013\\ \hline
Semester & Ganjil\\ \hline
Alokasi rencana & Remedial: 6 JP; pengayaan: 4 JP. Satu JP direncanakan 45 menit; penjadwalan mengikuti kebutuhan murid dan program madrasah.\\ \hline
\end{longtable}}\medskip
\vspace{8mm}\noindent\begin{minipage}[t]{.61\linewidth}Mengetahui,\\Kepala Madrasah,\\[22mm]Dr. Burhanuddin, S.Ag., M.Fil.I.\\NIP 197609152007101001\end{minipage}\hfill\begin{minipage}[t]{.35\linewidth}Gowa, 13 Juli 2026\\Guru Mata Pelajaran,\\[22mm]\textbf{Imran, S.Si.}\\NIP 199309082019031013\end{minipage}
\clearpage
\begin{center}\bfseries\large A. LANDASAN DAN ARAH PROGRAM\end{center}\medskip
{\small\setlength{\tabcolsep}{2mm}\renewcommand{\arraystretch}{1.35}\begin{longtable}{|>{\raggedright\arraybackslash}p{60.554mm}|>{\raggedright\arraybackslash}p{108.996mm}|}\hline
\rowcolor{gray!10}\textbf{Acuan} & \textbf{Penerapan}\\ \hline
\endfirsthead
\rowcolor{gray!10}\textbf{Acuan} & \textbf{Penerapan}\\ \hline
\endhead
KMA 450 Tahun 2024 sebagaimana diubah KMA 1503 Tahun 2025 & Layanan sesuai kebutuhan murid dan budaya cinta ilmu, kepedulian, serta saling membantu. [1]\\ \hline
Permendikdasmen 13 Tahun 2025; 1 Tahun 2026 & Kerangka kurikulum dan standar proses: perencanaan, pelaksanaan, serta penilaian proses. [2][3]\\ \hline
Permendikbudristek 21 Tahun 2022 & Asesmen berdasarkan bukti untuk menilai kemajuan dan tindak lanjut. [4]\\ \hline
Keputusan Kepala BSKAP 046/H/KR/2025 & CP Matematika Fase F, elemen 6.1 Bilangan. [5]\\ \hline
Panduan Pembelajaran dan Asesmen, revisi 2025 & Pendampingan untuk kompetensi yang belum dikuasai, asesmen formatif, dan pengayaan. [6]\\ \hline
\end{longtable}}\medskip
\par\medskip\noindent\textbf{Ringkasan CP yang relevan}\par\smallskip
CP Bilangan Fase F mencakup penalaran barisan dan deret aritmetika/geometri, penggunaannya pada bunga sederhana dan majemuk, pemodelan pinjaman serta investasi dengan anuitas, dan penyelidikan pengaruh parameter keuangan secara numerik atau grafis. [5]\par\medskip
{\small Kemendikdasmen pada 23 Juni 2026 menjelaskan bahwa pembaruan CP tahun 2026 terbatas pada Agama dan Budi Pekerti. CP Matematika tetap menggunakan keputusan tahun 2025. [7]\par}\medskip
\par\medskip\noindent\textbf{Tujuan dan lingkup program}\par\smallskip
Barisan dan deret aritmetika/geometri; bunga tunggal/majemuk; pengantar anuitas dan sensitivitas parameter.\par\medskip
{\small Program ditujukan untuk pemulihan kompetensi yang belum tercapai serta pendalaman bagi murid yang sudah menguasainya. Penempatan di kelas XI diselaraskan dengan ATP madrasah; CP merupakan target akhir fase. Tanggal pelaksanaan, peserta, dan hasil program diisi setelah kegiatan berdasarkan bukti nyata.\par}\medskip
\clearpage
\begin{center}\bfseries\large B. TUJUAN PEMBELAJARAN DAN KKTP\end{center}\medskip
{\small\setlength{\tabcolsep}{2mm}\renewcommand{\arraystretch}{1.35}\begin{longtable}{|>{\raggedright\arraybackslash}p{16.269mm}|>{\raggedright\arraybackslash}p{67.787mm}|>{\raggedright\arraybackslash}p{81.344mm}|}\hline
\rowcolor{gray!10}\textbf{Kode} & \textbf{Tujuan pembelajaran} & \textbf{Bukti esensial ketercapaian}\\ \hline
\endfirsthead
\rowcolor{gray!10}\textbf{Kode} & \textbf{Tujuan pembelajaran} & \textbf{Bukti esensial ketercapaian}\\ \hline
\endhead
TP1 & Menentukan pola dan suku ke-n barisan aritmetika. & Menetapkan suku pertama, beda, indeks awal, dan rumus; memeriksa hasil pada suku yang diketahui.\\ \hline
TP2 & Menentukan suku dan jumlah deret geometri berhingga. & Menetapkan rasio, menggunakan pangkat dan rumus jumlah dengan benar, serta memeriksa kasus sederhana.\\ \hline
TP3 & Memodelkan jumlah aritmetika/geometri dalam konteks. & Membedakan suku dan jumlah, memilih model, menentukan batas banyak suku, dan menafsirkan satuan.\\ \hline
TP4 & Membandingkan bunga tunggal dan majemuk. & Menggunakan bunga dan periode yang sepadan, menghitung saldo, serta menjelaskan asumsi model.\\ \hline
TP5 & Menganalisis pinjaman beranuitas dan perubahan parameter. & Menghitung pembayaran tetap akhir periode, memisahkan bunga/angsuran pokok, memeriksa sisa pinjaman, dan menafsirkan satu perubahan parameter.\\ \hline
\end{longtable}}\medskip
\par\medskip\noindent\textbf{Rubrik ketercapaian per TP}\par\smallskip
{\small\setlength{\tabcolsep}{2mm}\renewcommand{\arraystretch}{1.35}\begin{longtable}{|>{\raggedright\arraybackslash}p{41.693mm}|>{\raggedright\arraybackslash}p{127.857mm}|}\hline
\rowcolor{gray!10}\textbf{Level} & \textbf{Deskripsi operasional}\\ \hline
\endfirsthead
\rowcolor{gray!10}\textbf{Level} & \textbf{Deskripsi operasional}\\ \hline
\endhead
1 - Perlu dukungan intensif & Konsep utama belum tampak walau diberi pertanyaan penuntun.\\ \hline
2 - Berkembang & Sebagian langkah benar; masih ada kesalahan konsep atau arahan pada langkah utama.\\ \hline
3 - Tercapai & Seluruh bukti esensial TP tampak mandiri dengan hasil dan alasan yang dapat diperiksa.\\ \hline
4 - Melampaui & Memenuhi level 3 serta memberi pemeriksaan berbeda, generalisasi, atau argumen baru yang sah.\\ \hline
\end{longtable}}\medskip
{\small KKTP program: setiap TP minimal level 3. TP level 1-2 diberi remedial terarah; TP yang telah tercapai dapat mendapat pengayaan. Kriteria dan level ini ditetapkan guru untuk unit, bukan angka ketuntasan nasional. Kompetensi yang belum terbukti tidak ditutupi oleh kelebihan skor TP lain.\par}\medskip
\clearpage
\begin{center}\bfseries\large C. PETA KONSEP DAN NOTASI\end{center}\medskip
{\small\setlength{\tabcolsep}{2mm}\renewcommand{\arraystretch}{1.35}\begin{longtable}{|>{\raggedright\arraybackslash}p{29.188mm}|>{\raggedright\arraybackslash}p{58.376mm}|>{\raggedright\arraybackslash}p{77.835mm}|}\hline
\rowcolor{gray!10}\textbf{Pokok konsep} & \textbf{Hubungan / rumus} & \textbf{Syarat dan pemaknaan}\\ \hline
\endfirsthead
\rowcolor{gray!10}\textbf{Pokok konsep} & \textbf{Hubungan / rumus} & \textbf{Syarat dan pemaknaan}\\ \hline
\endhead
Aritmetika & \(U_n=a+(n-1)b\) & \(S_n=\frac n2\bigl(2a+(n-1)b\bigr)\).\\ \hline
Geometri & \(U_n=ar^{n-1}\) & \(S_n=a\frac{1-r^n}{1-r}\) untuk \(r\ne1\); jika \(r=1\), \(S_n=na\).\\ \hline
Bunga tunggal & \(F_n=P(1+ni)\) & Bunga dihitung dari modal awal; periode dan satuan bunga harus sama.\\ \hline
Bunga majemuk & \(F_n=P(1+i)^n\) & Bunga tiap periode ditambahkan ke saldo sebelum periode berikutnya.\\ \hline
Anuitas akhir periode & \(A=P\frac{i}{1-(1+i)^{-n}}\) & \(i>0\), \(n\) bulat positif; jika \(i=0\), \(A=P/n\).\\ \hline
Sisa pinjaman & \(B_k=(1+i)B_{k-1}-A\) & \(B_0=P\); bunga \(I_k=iB_{k-1}\); pokok \(C_k=A-I_k\).\\ \hline
\end{longtable}}\medskip
{\small \(U_n\): suku ke-\(n\); \(S_n\): jumlah \(n\) suku; \(a\): suku pertama; \(b\): beda; \(r\): rasio. Untuk barisan, indeks dimulai \(n=1\). Pada keuangan, \(P\): modal/pinjaman awal; \(i\): bunga per periode; \(F_n\): saldo akhir; \(A\): pembayaran tetap; \(B_k\): sisa pokok setelah pembayaran ke-\(k\).\par}\medskip
\par\medskip\noindent\textbf{Urutan pemeriksaan model}\par\smallskip
Kenali pola; tentukan indeks dan satuan; pilih rumus suku/jumlah; hitung; periksa kasus kecil; tafsirkan hasil dan batas model.\par\medskip
{\small Pada soal keuangan, seluruh angka merupakan simulasi pembelajaran. Bunga selalu dinyatakan per periode yang sama dengan pembayaran; pembayaran anuitas dilakukan di akhir periode, tanpa biaya tambahan. Pembulatan dicatat dan sisa pembayaran terakhir disesuaikan bila perlu.\par}\medskip
\clearpage
\begin{center}\bfseries\large LAMPIRAN 1. ASESMEN AWAL\end{center}\medskip
{\small Nama: \_\_\_\_\_\_\_\_\_\_\_\_\_\_\_\_\_\_\_\_\_\_\_\_\_\_\_\_\_\_\_\_\_\_  Kelas: \_\_\_\_\_\_\_\_\_\_\par}\medskip
{\small Waktu rencana: 30 menit. Kerjakan mandiri, sertakan langkah dan alasan; gunakan lembar tambahan atau sketsa bila perlu.\par}\medskip
\noindent\textbf{D1 / TP1.} Diketahui barisan \(3,7,11,15,\ldots\). Tentukan \(a\), \(b\), rumus \(U_n\), dan \(U_{10}\). Periksa rumus untuk \(n=1\) dan \(n=3\).\par\vspace{10mm}
\noindent\textbf{D2 / TP2.} Diketahui barisan \(2,6,18,\ldots\). Tentukan rasio, \(U_5\), dan \(S_5\). Jelaskan perbedaan suku ke-5 dan jumlah lima suku.\par\vspace{10mm}
\noindent\textbf{D3 / TP3.} Baris pertama kursi berisi 8 kursi dan setiap baris berikutnya bertambah 2 kursi. Ada 12 baris. Tentukan banyak kursi baris terakhir dan jumlah seluruh kursi. Tuliskan asumsi.\par\vspace{10mm}
\noindent\textbf{D4 / TP4.} Modal Rp1.000.000 mendapat bunga 10\% per tahun selama dua tahun. Bandingkan saldo bunga tunggal dan majemuk tahunan tanpa biaya atau penarikan. Jelaskan penyebab selisihnya.\par\vspace{10mm}
\noindent\textbf{D5 / TP5.} Pinjaman Rp210.000 dibayar dengan dua pembayaran tetap pada akhir setiap periode. Bunga simulasi 10\% per periode dari sisa pokok. Hitung anuitas dan susun dua baris tabel bunga, angsuran pokok, dan sisa pinjaman.\par\vspace{10mm}
{\small Guru mengonfirmasi pemahaman lewat pertanyaan lisan bila bukti tertulis belum cukup. Kelompok layanan dibentuk per TP dan dapat berubah mengikuti kemajuan.\par}\medskip
\clearpage
\begin{center}\bfseries\large D. RENCANA PROGRAM REMEDIAL\end{center}\medskip
Peserta adalah murid dengan bukti TP level 1-2. Satu siklus direncanakan 3 pertemuan × 2 JP, hanya untuk kompetensi yang masih memerlukan bantuan. Durasi disesuaikan perkembangan peserta.\par\medskip
{\small\setlength{\tabcolsep}{2mm}\renewcommand{\arraystretch}{1.35}\begin{longtable}{|>{\raggedright\arraybackslash}p{16.847mm}|>{\raggedright\arraybackslash}p{89.049mm}|>{\raggedright\arraybackslash}p{16.847mm}|>{\raggedright\arraybackslash}p{38.507mm}|}\hline
\rowcolor{gray!10}\textbf{Sesi} & \textbf{Kegiatan pendampingan} & \textbf{Waktu} & \textbf{Bukti kemajuan}\\ \hline
\endfirsthead
\rowcolor{gray!10}\textbf{Sesi} & \textbf{Kegiatan pendampingan} & \textbf{Waktu} & \textbf{Bukti kemajuan}\\ \hline
\endhead
R1\newline TP1-TP2 & Memahami: wawancara kesalahan (15 menit). Mengaplikasi: contoh terurai dan latihan terbimbing (35 menit), latihan mandiri (30 menit). Merefleksi: tiket keluar (10 menit). Kartu beda/rasio, daftar suku, garis indeks, serta perbedaan suku dan jumlah. & 2 JP\newline 90 menit & Catat alasan, hasil, bantuan yang masih diperlukan, dan level TP terbaru.\\ \hline
R2\newline TP3-TP4 & Memahami: wawancara kesalahan (15 menit). Mengaplikasi: contoh terurai dan latihan terbimbing (35 menit), latihan mandiri (30 menit). Merefleksi: tiket keluar (10 menit). Model kursi/tanaman, tabel akumulasi, dan saldo bunga tahun demi tahun. & 2 JP\newline 90 menit & Catat alasan, hasil, bantuan yang masih diperlukan, dan level TP terbaru.\\ \hline
R3\newline TP5 & Memahami: identifikasi model (15 menit). Mengaplikasi: latihan terarah (20 menit), asesmen ulang TP sasaran (40 menit). Merefleksi: umpan balik dan tindak lanjut (15 menit). Tabel sisa pokok, bunga, dan pembayaran tetap; satu parameter diubah setiap kali. & 2 JP\newline 90 menit & Catat alasan, hasil, bantuan yang masih diperlukan, dan level TP terbaru.\\ \hline
\end{longtable}}\medskip
\par\medskip\noindent\textbf{Dukungan dan nilai pembelajaran}\par\smallskip
Gunakan tabel, sketsa, kartu konsep, serta contoh konkret. Guru memodelkan cara berpikir lalu mengurangi bantuan. Tutor sebaya mendukung diskusi dan tidak menggantikan pendampingan guru. Sediakan lembar cetak, instruksi yang mudah dibaca, tambahan waktu atau jawaban lisan sesuai kebutuhan.\par\medskip
{\small Cinta ilmu diwujudkan melalui ketelitian dan kejujuran pemeriksaan; cinta diri dan sesama melalui bantuan tanpa mempermalukan kesalahan. Konteks tanaman atau taman menumbuhkan kepedulian lingkungan.\par}\medskip
\par\medskip\noindent\textbf{Keputusan setelah layanan}\par\smallskip
{\small TP level 3 kembali ke kegiatan reguler atau mendapat pengayaan. Jika belum tercapai, ubah bantuan dan jadwalkan siklus singkat berikutnya. Simpan bukti awal, lembar kerja, asesmen ulang, dan umpan balik. Pelaporan nilai mengikuti kebijakan madrasah tanpa pembatasan otomatis nilai remedial dalam program ini.\par}\medskip
\clearpage
\begin{center}\bfseries\large LAMPIRAN 2. LEMBAR KERJA REMEDIAL\end{center}\medskip
{\small Nama: \_\_\_\_\_\_\_\_\_\_\_\_\_\_\_\_\_\_\_\_\_\_\_\_\_\_\_\_\_\_\_\_\_\_  Kelas: \_\_\_\_\_\_\_\_\_\_\par}\medskip
{\small Guru memilih butir sesuai TP sasaran. Jelaskan alasan pemilihan rumus dan tunjukkan pemeriksaan hasil.\par}\medskip
\noindent\textbf{R1 / TP1.} Barisan \(5,8,11,\ldots\). Tentukan beda, rumus \(U_n\), \(U_8\), dan jumlah delapan sukunya. Bantuan: tandai bahwa suku pertama bersesuaian dengan \(n=1\).\par\vspace{10mm}
\noindent\textbf{R2 / TP2.} Barisan \(3,6,12,\ldots\). Tentukan \(r\), \(U_6\), dan \(S_6\). Periksa jumlah tiga suku pertama secara langsung sebelum menggunakan rumus.\par\vspace{10mm}
\noindent\textbf{R3 / TP3.} Terdapat 10 baris tanaman. Baris pertama berisi 10 tanaman, lalu bertambah 3 tanaman tiap baris. Hitung baris terakhir dan jumlah semua tanaman. Jelaskan alasan memilih deret aritmetika.\par\vspace{10mm}
\noindent\textbf{R4 / TP4.} Bandingkan saldo modal Rp2.000.000 selama dua tahun pada bunga tunggal 5\% per tahun dan bunga majemuk 5\% dengan kapitalisasi tahunan. Buat tabel tahun 0, 1, dan 2.\par\vspace{10mm}
\noindent\textbf{R5 / TP5.} Pinjaman Rp420.000 dibayar dua kali pada akhir periode, bunga simulasi 10\% tiap periode. Gunakan rumus anuitas; pisahkan bunga dan pokok pada tiap pembayaran, lalu periksa sisa akhir nol.\par\vspace{10mm}
{\small Kesalahan yang saya perbaiki: \_\_\_\_\_\_\_\_\_\_\_\_\_\_\_\_\_\_\_\_\_\_\_\_\_\_\_\_\_\_\_\_\_\_\_\_\_\_\_\_\newline Cara saya memeriksa jawaban: \_\_\_\_\_\_\_\_\_\_\_\_\_\_\_\_\_\_\_\_\_\_\_\_\_\_\_\_\_\_\_\_\_\_\_\_\_\_\_\_\_\par}\medskip
\clearpage
\begin{center}\bfseries\large LAMPIRAN 3. ASESMEN ULANG\end{center}\medskip
{\small Nama: \_\_\_\_\_\_\_\_\_\_\_\_\_\_\_\_\_\_\_\_\_\_\_\_\_\_\_\_\_\_\_\_\_\_  Kelas: \_\_\_\_\_\_\_\_\_\_\par}\medskip
{\small Waktu rencana: 40 menit. Kerjakan mandiri; asesmen dipilih sesuai TP sasaran remedial.\par}\medskip
\noindent\textbf{U1 / TP1.} Barisan \(4,9,14,\ldots\). Tentukan \(a\), \(b\), rumus \(U_n\), dan \(U_{12}\); periksa dua suku awal.\par\vspace{10mm}
\noindent\textbf{U2 / TP2.} Barisan \(4,8,16,\ldots\). Tentukan rasio, \(U_6\), dan \(S_6\). Jelaskan perbedaan hasil kedua besaran itu.\par\vspace{10mm}
\noindent\textbf{U3 / TP3.} Sebuah susunan kursi memiliki 10 baris. Baris pertama berisi 12 kursi dan tiap baris bertambah 2 kursi. Hitung banyak kursi pada baris terakhir serta jumlah keseluruhannya.\par\vspace{10mm}
\noindent\textbf{U4 / TP4.} Modal Rp1.000.000 diberi bunga 20\% per tahun selama dua tahun. Bandingkan bunga tunggal dengan bunga majemuk tahunan dan jelaskan satu asumsi model.\par\vspace{10mm}
\noindent\textbf{U5 / TP5.} Pinjaman Rp630.000 dibayar dengan dua anuitas akhir periode, bunga simulasi 10\% tiap periode dari sisa pokok. Tentukan pembayaran tetap dan lengkapi tabel bunga, pokok, serta sisa pinjaman.\par\vspace{10mm}
{\small Penetapan hasil memakai rubrik bagian B. TP tercapai ketika seluruh bukti esensialnya tampak mandiri minimal level 3. Periksa konsep, alasan, satuan, dan kewajaran hasil.\par}\medskip
\clearpage
\begin{center}\bfseries\large E. RENCANA PROGRAM PENGAYAAN\end{center}\medskip
Peserta telah mencapai minimal level 3 pada TP sasaran. Pilih satu proyek; peserta tetap menjelaskan kontribusi dan pemikirannya sendiri. Pengayaan memperluas pemodelan dan penalaran, tidak sekadar menambah soal sejenis.\par\medskip
{\small\setlength{\tabcolsep}{2mm}\renewcommand{\arraystretch}{1.35}\begin{longtable}{|>{\raggedright\arraybackslash}p{15.752mm}|>{\raggedright\arraybackslash}p{112.892mm}|>{\raggedright\arraybackslash}p{36.756mm}|}\hline
\rowcolor{gray!10}\textbf{Sesi} & \textbf{Aktivitas} & \textbf{Waktu / produk}\\ \hline
\endfirsthead
\rowcolor{gray!10}\textbf{Sesi} & \textbf{Aktivitas} & \textbf{Waktu / produk}\\ \hline
\endhead
P1 & Pilih konteks dan nyatakan asumsi (15 menit); eksplorasi rumus/tabel/sketsa (60 menit); pemeriksaan dan refleksi (15 menit). & 2 JP / 90 menit\newline Draf solusi dan bukti pemeriksaan.\\ \hline
P2 & Perbaiki model berdasar umpan balik (35 menit); presentasi serta verifikasi individual (35 menit); evaluasi keterbatasan (20 menit). & 2 JP / 90 menit\newline Laporan 1-2 halaman dan refleksi individual.\\ \hline
\end{longtable}}\medskip
\par\medskip\noindent\textbf{Proyek A - Membandingkan dua rencana penanaman}\par\smallskip
Rencana A menanam 6 bibit pada pekan pertama dan bertambah 2 bibit tiap pekan. Rencana B menanam 3 bibit pada pekan pertama dan berlipat dua tiap pekan. Buat rumus dan tabel enam pekan; bandingkan total bibit, tentukan pekan pertama jumlah mingguan B melampaui A, dan evaluasi batas model terhadap lahan serta ketersediaan bibit.\par\medskip
\par\medskip\noindent\textbf{Proyek B - Membandingkan skema pembayaran}\par\smallskip
Gunakan pinjaman simulasi Rp210.000 selama dua periode. Skema flat: pokok Rp105.000 tiap periode dan bunga 10\% dari pinjaman awal tiap periode. Skema anuitas: pembayaran akhir periode dengan bunga 10\% dari sisa pokok. Bandingkan pembayaran dan total bunga. Pada anuitas, selidiki dampak bunga menjadi 20\% dengan lama tetap dua periode, lalu lama menjadi tiga periode pada bunga 10\%. Nyatakan asumsi dan pembulatan.\par\medskip
{\small Produk memuat masalah, variabel/unsur, rumus, tabel atau sketsa, pemeriksaan, kesimpulan, dan batas model. Gunakan data soal sebagai simulasi; data nyata hanya dicantumkan bila benar-benar diperoleh.\par}\medskip
\clearpage
\begin{center}\bfseries\large F. RUBRIK PENGAYAAN DAN REFLEKSI\end{center}\medskip
{\small\setlength{\tabcolsep}{2mm}\renewcommand{\arraystretch}{1.35}\begin{longtable}{|>{\raggedright\arraybackslash}p{27.873mm}|>{\raggedright\arraybackslash}p{26.606mm}|>{\raggedright\arraybackslash}p{29.140mm}|>{\raggedright\arraybackslash}p{35.474mm}|>{\raggedright\arraybackslash}p{38.008mm}|}\hline
\rowcolor{gray!10}\textbf{Aspek} & \textbf{Level 1} & \textbf{Level 2} & \textbf{Level 3} & \textbf{Level 4}\\ \hline
\endfirsthead
\rowcolor{gray!10}\textbf{Aspek} & \textbf{Level 1} & \textbf{Level 2} & \textbf{Level 3} & \textbf{Level 4}\\ \hline
\endhead
Model / representasi & Belum sesuai konteks. & Sebagian sesuai; asumsi belum jelas. & Model dan representasi tepat serta beralasan. & Menguji batas atau memberi alternatif yang sah.\\ \hline
Prosedur / argumen & Langkah utama belum benar. & Memerlukan arahan pada langkah utama. & Perhitungan serta alasan runtut dan mandiri. & Memberi generalisasi atau pemeriksaan berbeda.\\ \hline
Evaluasi konteks & Belum menafsirkan hasil. & Tafsiran sebagian sesuai. & Kesimpulan wajar; satuan dan asumsi jelas. & Menguji keterbatasan dan memberi usul validasi.\\ \hline
Komunikasi individual & Belum menjelaskan kontribusi. & Bergantung pada pasangan. & Menjelaskan kontribusi dan alasan sendiri. & Menanggapi sanggahan serta memperbaiki penjelasan.\\ \hline
\end{longtable}}\medskip
{\small Target produk: minimal level 3 pada keempat aspek. Skala ditetapkan guru untuk tugas ini. Penilaian kelompok dilengkapi verifikasi individual; kerja sama diberi umpan balik dan tidak menggantikan bukti matematika.\par}\medskip
\par\medskip\noindent\textbf{Refleksi murid}\par\smallskip
\noindent\textbf{1.} Konsep apa yang membantu saya memilih rumus atau hubungan?\par\vspace{6mm}
\noindent\textbf{2.} Bagaimana saya memeriksa indeks, satuan, atau unsur pada sketsa?\par\vspace{6mm}
\noindent\textbf{3.} Asumsi apa yang membatasi penggunaan hasil?\par\vspace{6mm}
\noindent\textbf{4.} Bantuan dan langkah belajar berikutnya apa yang saya perlukan?\par\vspace{6mm}
\par\medskip\noindent\textbf{Refleksi guru}\par\smallskip
{\small Kesulitan yang masih muncul: \_\_\_\_\_\_\_\_\_\_\_\_\_\_\_\_\_\_\_\_\_\_\_\_\_\_\_\_\_\_\_\_\_\_\_\_\_\_\_\_\_\newline Bantuan yang efektif dan buktinya: \_\_\_\_\_\_\_\_\_\_\_\_\_\_\_\_\_\_\_\_\_\_\_\_\_\_\_\_\_\_\_\_\_\_\_\_\_\newline Perubahan untuk siklus berikutnya: \_\_\_\_\_\_\_\_\_\_\_\_\_\_\_\_\_\_\_\_\_\_\_\_\_\_\_\_\_\_\_\_\_\_\_\_\par}\medskip
\clearpage
\begin{center}\bfseries\large LAMPIRAN 4. KUNCI ASESMEN AWAL\\[2pt] DAN LEMBAR KERJA REMEDIAL\end{center}\medskip
{\small\setlength{\tabcolsep}{2mm}\renewcommand{\arraystretch}{1.35}\begin{longtable}{|>{\raggedright\arraybackslash}p{14.685mm}|>{\raggedright\arraybackslash}p{154.865mm}|}\hline
\rowcolor{gray!10}\textbf{Butir} & \textbf{Jawaban dan bukti yang diperiksa}\\ \hline
\endfirsthead
\rowcolor{gray!10}\textbf{Butir} & \textbf{Jawaban dan bukti yang diperiksa}\\ \hline
\endhead
D1 & \(a=3\), \(b=4\); \(U_n=4n-1\); \(U_{10}=39\). Pemeriksaan: \(U_1=3\), \(U_3=11\).\\ \hline
D2 & \(r=3\); \(U_5=2\cdot3^4=162\); \(S_5=2\frac{3^5-1}{3-1}=242\). Suku berbeda dengan jumlah.\\ \hline
D3 & \(U_{12}=8+11\cdot2=30\) kursi; \(S_{12}=\frac{12}{2}(8+30)=228\) kursi. Semua baris mengikuti pola tanpa pengurangan kursi.\\ \hline
D4 & Tunggal: \(1\,000\,000(1+2\cdot0{,}1)=1\,200\,000\). Majemuk: \(1\,000\,000(1{,}1)^2=1\,210\,000\). Selisih Rp10.000 berasal dari bunga atas bunga tahun pertama.\\ \hline
D5 & \(A=121\,000\). Periode 1: bunga Rp21.000, pokok Rp100.000, sisa Rp110.000. Periode 2: bunga Rp11.000, pokok Rp110.000, sisa nol.\\ \hline
R1 & \(b=3\); \(U_n=3n+2\); \(U_8=26\); \(S_8=\frac82(5+26)=124\).\\ \hline
R2 & \(r=2\); \(U_6=96\); \(S_6=3(2^6-1)=189\). Tiga suku pertama berjumlah 21.\\ \hline
R3 & \(U_{10}=10+9\cdot3=37\); \(S_{10}=\frac{10}{2}(10+37)=235\) tanaman.\\ \hline
R4 & Saldo tunggal tahun 0, 1, 2: Rp2.000.000; Rp2.100.000; Rp2.200.000. Majemuk: Rp2.000.000; Rp2.100.000; Rp2.205.000.\\ \hline
R5 & \(A=242\,000\). Periode 1: bunga Rp42.000, pokok Rp200.000, sisa Rp220.000. Periode 2: bunga Rp22.000, pokok Rp220.000, sisa nol.\\ \hline
\end{longtable}}\medskip
{\small Jawaban ekuivalen diterima jika langkah dan alasan sah. Kunci membantu pemeriksaan; level TP tetap ditentukan dengan memperhatikan seluruh bukti esensial dan kemandirian.\par}\medskip
\clearpage
\begin{center}\bfseries\large LAMPIRAN 5. KUNCI ASESMEN ULANG\\[2pt] DAN PENGAYAAN\end{center}\medskip
{\small\setlength{\tabcolsep}{2mm}\renewcommand{\arraystretch}{1.35}\begin{longtable}{|>{\raggedright\arraybackslash}p{31.791mm}|>{\raggedright\arraybackslash}p{137.759mm}|}\hline
\rowcolor{gray!10}\textbf{Butir} & \textbf{Jawaban / panduan pemeriksaan}\\ \hline
\endfirsthead
\rowcolor{gray!10}\textbf{Butir} & \textbf{Jawaban / panduan pemeriksaan}\\ \hline
\endhead
U1 & \(a=4\), \(b=5\); \(U_n=5n-1\); \(U_{12}=59\).\\ \hline
U2 & \(r=2\); \(U_6=128\); \(S_6=4(2^6-1)=252\).\\ \hline
U3 & \(U_{10}=30\) kursi; \(S_{10}=\frac{10}{2}(12+30)=210\) kursi.\\ \hline
U4 & Tunggal Rp1.400.000; majemuk Rp1.440.000; selisih Rp40.000. Bunga dan kapitalisasi tahunan; tanpa biaya atau penarikan.\\ \hline
U5 & \(A=363\,000\). Periode 1: bunga Rp63.000, pokok Rp300.000, sisa Rp330.000. Periode 2: bunga Rp33.000, pokok Rp330.000, sisa nol.\\ \hline
Proyek A & \(U_n^{(A)}=6+2(n-1)\), \(S_6^{(A)}=66\). \(U_n^{(B)}=3\cdot2^{n-1}\), \(S_6^{(B)}=189\). Jumlah mingguan B pertama melampaui A pada pekan ke-3: 12 dibanding 10.\\ \hline
Proyek B: dasar & Flat: pembayaran Rp126.000 per periode, total Rp252.000, bunga Rp42.000. Anuitas: Rp121.000 per periode, total Rp242.000, bunga Rp31.000. Perbedaan berasal dari dasar penghitungan bunga.\\ \hline
Proyek B: parameter & Bunga 20\%, dua periode: \(A\approx137\,454{,}55\). Bunga 10\%, tiga periode: \(A\approx84\,444{,}11\); total bunga sekitar Rp43.332,33. Periode lebih panjang menurunkan pembayaran per periode, tetapi dapat meningkatkan total bunga. Pembayaran terakhir disesuaikan jika ada pembulatan.\\ \hline
\end{longtable}}\medskip
\clearpage
\begin{center}\bfseries\large LAMPIRAN 6. FORMAT PELAKSANAAN\\[2pt] DAN TINDAK LANJUT\end{center}\medskip
\par\medskip\noindent\textbf{1. Pemetaan dan keputusan per peserta}\par\smallskip
{\small\setlength{\tabcolsep}{2mm}\renewcommand{\arraystretch}{1.35}\begin{longtable}{|>{\raggedright\arraybackslash}p{10.030mm}|>{\raggedright\arraybackslash}p{37.611mm}|>{\raggedright\arraybackslash}p{20.059mm}|>{\raggedright\arraybackslash}p{17.552mm}|>{\raggedright\arraybackslash}p{22.566mm}|>{\raggedright\arraybackslash}p{45.133mm}|}\hline
\rowcolor{gray!10}\textbf{No} & \textbf{Nama / kelas} & \textbf{TP sasaran} & \textbf{Level awal} & \textbf{Layanan} & \textbf{Level akhir / tindak lanjut}\\ \hline
\endfirsthead
\rowcolor{gray!10}\textbf{No} & \textbf{Nama / kelas} & \textbf{TP sasaran} & \textbf{Level awal} & \textbf{Layanan} & \textbf{Level akhir / tindak lanjut}\\ \hline
\endhead
1 & \rule{0pt}{8mm} &  &  &  & \\ \hline
2 & \rule{0pt}{8mm} &  &  &  & \\ \hline
3 & \rule{0pt}{8mm} &  &  &  & \\ \hline
4 & \rule{0pt}{8mm} &  &  &  & \\ \hline
5 & \rule{0pt}{8mm} &  &  &  & \\ \hline
6 & \rule{0pt}{8mm} &  &  &  & \\ \hline
\end{longtable}}\medskip
\par\medskip\noindent\textbf{2. Jurnal kegiatan}\par\smallskip
{\small\setlength{\tabcolsep}{2mm}\renewcommand{\arraystretch}{1.35}\begin{longtable}{|>{\raggedright\arraybackslash}p{26.434mm}|>{\raggedright\arraybackslash}p{29.078mm}|>{\raggedright\arraybackslash}p{47.582mm}|>{\raggedright\arraybackslash}p{58.156mm}|}\hline
\rowcolor{gray!10}\textbf{Tanggal / JP} & \textbf{TP / kelompok} & \textbf{Bantuan atau tugas} & \textbf{Bukti belajar / umpan balik}\\ \hline
\endfirsthead
\rowcolor{gray!10}\textbf{Tanggal / JP} & \textbf{TP / kelompok} & \textbf{Bantuan atau tugas} & \textbf{Bukti belajar / umpan balik}\\ \hline
\endhead
\rule{0pt}{8mm} &  &  & \\ \hline
\rule{0pt}{8mm} &  &  & \\ \hline
\rule{0pt}{8mm} &  &  & \\ \hline
\end{longtable}}\medskip
\par\medskip\noindent\textbf{3. Evaluasi program}\par\smallskip
{\small\setlength{\tabcolsep}{2mm}\renewcommand{\arraystretch}{1.35}\begin{longtable}{|>{\raggedright\arraybackslash}p{77.826mm}|>{\raggedright\arraybackslash}p{91.724mm}|}\hline
Jumlah peserta remedial / pengayaan & \_\_\_\_\_\_\_\_\_\_\_\_ / \_\_\_\_\_\_\_\_\_\_\_\_\\ \hline
TP tercapai setelah layanan & \_\_\_\_\_\_\_\_\_\_\_\_\_\_\_\_\_\_\_\_\_\_\_\_\_\_\_\_\_\_\_\_\\ \hline
TP masih perlu dukungan & \_\_\_\_\_\_\_\_\_\_\_\_\_\_\_\_\_\_\_\_\_\_\_\_\_\_\_\_\_\_\_\_\\ \hline
Tindak lanjut dan jadwal berikutnya & \_\_\_\_\_\_\_\_\_\_\_\_\_\_\_\_\_\_\_\_\_\_\_\_\_\_\_\_\_\_\_\_\\ \hline
Dokumen bukti & Asesmen awal, lembar kerja, asesmen ulang, produk, rubrik, dan umpan balik.\\ \hline
\end{longtable}}\medskip
{\small Diisi setelah kegiatan berdasarkan bukti nyata; kehadiran atau selesainya tugas tidak menggantikan bukti ketercapaian kompetensi.\par}\medskip
\clearpage
\begin{center}\bfseries\large DAFTAR RUJUKAN RESMI\end{center}\medskip
[1] Kementerian Agama. KMA Nomor 1503 Tahun 2025 tentang perubahan atas KMA Nomor 450 Tahun 2024; terutama Bab IV Pembelajaran dan Asesmen.\par\medskip
\href{https://sulteng.kemenag.go.id/files/show/20251020162801surat-penyampaian-kma-kurikulum-2025%2Blampiran_compressed-tzR.pdf}{\small\textcolor{blue}{Buka sumber resmi}}\par
\vspace{8.50393701pt}
[2] Permendikdasmen Nomor 13 Tahun 2025 tentang perubahan ketentuan kurikulum; tercantum dalam dasar hukum Keputusan Kepala BSKAP 046/H/KR/2025.\par\medskip
\href{https://bpmpbabel.kemendikdasmen.go.id/PPID/wp-content/uploads/2025/07/KepKaBSKAP-046_2025-ttg-CP.pdf}{\small\textcolor{blue}{Buka sumber resmi}}\par
\vspace{8.50393701pt}
[3] Kemendikdasmen. Abstrak resmi Permendikdasmen Nomor 1 Tahun 2026 tentang Standar Proses. Berlaku 5 Januari 2026 dan menggantikan standar proses tahun 2022.\par\medskip
\href{https://jdih.kemendikdasmen.go.id/sjdih/siperpu/dokumen/abstraksi/Abstrak%20RPM%20Standar%20Proses.pdf}{\small\textcolor{blue}{Buka sumber resmi}}\par
\vspace{8.50393701pt}
[4] Pusat Standar dan Kebijakan Pendidikan. Laman Standar Nasional Pendidikan, bagian Standar Penilaian; memuat rujukan Permendikbudristek Nomor 21 Tahun 2022.\par\medskip
\href{https://pskp.kemendikdasmen.go.id/standar-nasional-pendidikan-a/standar-proses}{\small\textcolor{blue}{Buka sumber resmi}}\par
\vspace{8.50393701pt}
[5] Keputusan Kepala BSKAP Nomor 046/H/KR/2025. Lampiran II, Matematika, Fase F, elemen 6.1 Bilangan.\par\medskip
\href{https://bpmpbabel.kemendikdasmen.go.id/PPID/wp-content/uploads/2025/07/KepKaBSKAP-046_2025-ttg-CP.pdf}{\small\textcolor{blue}{Buka sumber resmi}}\par
\vspace{8.50393701pt}
[6] Kemendikdasmen. Panduan Pembelajaran dan Asesmen, Edisi Revisi 2025. Bagian kriteria ketercapaian, remedial, pengayaan, dan pengolahan hasil asesmen.\par\medskip
\href{https://kurikulum.kemendikdasmen.go.id/file/1755668120_manage_file.pdf}{\small\textcolor{blue}{Buka sumber resmi}}\par
\vspace{8.50393701pt}
[7] Kemendikdasmen. Siaran pers 23 Juni 2026 tentang ruang lingkup pembaruan CP melalui Keputusan Kepala BKPDM Nomor 020 Tahun 2026.\par\medskip
\href{https://www.kemendikdasmen.go.id/siaran-pers/15636-capaian-pembelajaran-baru-telah-terbit-yang-berubah-hanya-mata-pelajaran-agama-dan-budi-pekerti}{\small\textcolor{blue}{Buka sumber resmi}}\par
\vspace{8.50393701pt}
{\small Pemeriksaan sumber: 31 Oktober 2026. Rencana waktu, perangkat tugas, rubrik, dan format administrasi dalam program ini disusun untuk kebutuhan unit oleh guru dan dapat diselaraskan dengan ATP serta kebijakan penilaian madrasah.\par}\medskip
\end{document}
